\chapter{\ifproject%
\ifenglish Experimentation and Results\else การทดลองและผลลัพธ์\fi
\else%
\ifenglish System Evaluation\else การประเมินระบบ\fi
\fi}

ในบทนี้นำเสนอการทดลองเบื้องต้นที่ดำเนินการแล้ว ได้แก่
(1) การทดลองส่ง Raw Logs ให้ LLM เชิงพาณิชย์วิเคราะห์โดยตรง เพื่อยืนยันปัญหาของโครงงาน
(2) การทดลองเพื่อตอบคำถามว่า LLM ขนาดเล็กยังจำเป็นต้องใช้ RAG หรือไม่
และ (3) แผนการทดลองหลักและการประเมินผล

\section{การทดลองเบื้องต้น: ส่ง Raw Logs ให้ LLM วิเคราะห์โดยตรง}

\subsection{วิธีการทดลอง}
ดึง Windows Sysmon Logs ของชุดข้อมูล T1087.001 จาก Splunk ผ่าน API ได้เป็นไฟล์ JSON
แล้วอัปโหลดไฟล์ให้ LLM เชิงพาณิชย์ 3 โมเดลวิเคราะห์โดยตรง
ด้วยคำถามว่า ``นี่คือ Attack Technique อะไร เพราะอะไรถึงคิดแบบนั้น''
คำตอบที่คาดหวังคือ Account Discovery: Local Account (T1087.001)

\subsection{ผลการทดลอง}
\begin{table}[ht]
  \centering
  \caption{ผลการส่ง Raw Logs ของ T1087.001 ให้ LLM เชิงพาณิชย์วิเคราะห์โดยตรง}
  \label{tab:raw-llm}
  \begin{tabular}{|L{0.25\textwidth}|L{0.45\textwidth}|L{0.17\textwidth}|}
    \hline
    \textbf{โมเดล} & \textbf{คำตอบ} & \textbf{ผล} \\ \hline
    GPT-5.6 Luna      & T1003 OS Credential Dumping & ผิด \\ \hline
    Gemini 3.1 Pro    & ไม่พบพฤติกรรมการโจมตีที่ชัดเจน (Logs เป็นเพียงการทำงานของ Splunk Agent) & ผิด \\ \hline
    Claude Sonnet 5 (High) & T1087 Account Discovery & ถูกระดับ Technique \\ \hline
  \end{tabular}
\end{table}

จากตารางที่~\ref{tab:raw-llm} พบว่าแม้เป็นโมเดลขนาดใหญ่ แต่เมื่อได้รับ Raw Logs ที่มี Noise จำนวนมาก
บางโมเดลถูกดึงความสนใจไปที่ Events ที่ไม่เกี่ยวข้อง (เช่น การเข้าถึง Process ของ Splunk Agent)
และไม่มีโมเดลใดระบุถึงระดับ Sub-technique ได้ถูกต้อง
ผลนี้สนับสนุนสมมติฐานว่าการสร้าง Investigation Context ที่คัดกรองแล้วมีความจำเป็น

\section{การทดลองนำร่อง: ความจำเป็นของ RAG สำหรับ LLM ขนาดเล็ก}

\subsection{วัตถุประสงค์และการตั้งค่า}
เพื่อทดสอบว่า เมื่อผ่าน Event Selection แล้ว LLM ขนาดเล็ก (Qwen3.5:9B) ยังต้องใช้ RAG หรือไม่
และ Query แบบใดเหมาะกับการค้นหาในฐานความรู้
Input ของการทดลองคือผลลัพธ์จาก Event Selection แบบ Top-$k$ ของชุดข้อมูล T1087.001
ประกอบด้วย Alert และ 16 Events (EID 1 จำนวน 3 Events ได้แก่ \texttt{cmd.exe} $\rightarrow$ \texttt{net users}
$\rightarrow$ \texttt{net1 users} และที่เหลือเป็น EID 10) ในรูปแบบ JSON
เงื่อนไขการทดลองแสดงในตารางที่~\ref{tab:rag-conditions}

\begin{table}[ht]
  \centering
  \caption{เงื่อนไขการทดลองความจำเป็นของ RAG}
  \label{tab:rag-conditions}
  \begin{tabular}{|L{0.25\textwidth}|c|L{0.52\textwidth}|}
    \hline
    \textbf{การทดลอง} & \textbf{RAG} & \textbf{คำอธิบาย} \\ \hline
    1. Direct-Log      & ไม่ใช้ & ส่ง Representation Log ให้ LLM โดยตรง \\ \hline
    2. Direct-Summary  & ไม่ใช้ & ให้ LLM แปลง Log เป็นคำอธิบายพฤติกรรมภาษาคนก่อน แล้วจึงให้ระบุเทคนิค \\ \hline
    3.1 RAG-Log ($k=10$) & ใช้ & ใช้ Representation Log เป็น Query ค้นหา RAG ได้ Top-10 ให้ LLM \\ \hline
    3.2 RAG-Log ($k=5$)  & ใช้ & ใช้ Representation Log เป็น Query ค้นหา RAG ได้ Top-5 ให้ LLM \\ \hline
    4. RAG-Summary     & ใช้ & ใช้คำอธิบายพฤติกรรมภาษาคน (จากการทดลองที่ 2) เป็น Query ค้นหา RAG ได้ Top-$k$ ให้ LLM \\ \hline
  \end{tabular}
\end{table}

\subsection{ผลการทดลอง}
ผลการทดลองทั้ง 5 เงื่อนไข ได้แก่ อันดับของเฉลยในผลลัพธ์ของ RAG และคำตอบของ LLM แสดงในตารางที่~\ref{tab:rag-results}

\begin{table}[ht]
  \centering
  \caption{ผลการทดลองความจำเป็นของ RAG (คำตอบที่ถูกต้องคือ T1087.001)}
  \label{tab:rag-results}
  \small
  \begin{tabular}{|l|c|l|c|}
    \hline
    \textbf{การทดลอง} & \textbf{อันดับเฉลยใน RAG} & \textbf{คำตอบของ LLM} & \textbf{ผล} \\ \hline
    1. Direct-Log        & --                & T1083 File and Directory Discovery & ผิด \\ \hline
    2. Direct-Summary    & --                & T1078 Valid Accounts                & ผิด \\ \hline
    3.1 RAG-Log ($k=10$) & อันดับ 6          & T1087.001                           & ถูก \\ \hline
    3.2 RAG-Log ($k=5$)  & ไม่อยู่ใน Top-5   & T1033 System Owner/User Discovery   & ผิด \\ \hline
    4. RAG-Summary       & อันดับ 1          & T1087.001                           & ถูก \\ \hline
  \end{tabular}
\end{table}

\subsection{การวิเคราะห์ผล}
จากตารางที่~\ref{tab:rag-results} พบว่า
\begin{enumerate}
  \item เมื่อไม่ใช้ RAG (การทดลองที่ 1 และ 2) LLM ขนาดเล็กตอบผิดทั้งสองกรณี
        แม้ Input จะผ่านการคัดเลือก Event มาแล้ว แสดงว่าโมเดลขนาด 9B ยังขาดความรู้เฉพาะด้าน ATT\&CK
        ที่จะแยกแยะเทคนิคที่ใกล้เคียงกันได้
  \item LLM ตอบถูกเฉพาะเมื่อมี T1087.001 ซึ่งเป็นเฉลยอยู่ใน Top-$k$ ที่ได้จาก RAG
        ดังนั้นคุณภาพของ Retrieval เป็นปัจจัยสำคัญต่อความแม่นยำ
  \item การใช้คำอธิบายพฤติกรรมภาษาคนเป็น Query (RAG-Summary) ทำให้เฉลยขึ้นมาอยู่อันดับ 1
        ดีกว่าการใช้ JSON โดยตรง (อันดับ 6) เนื่องจากคำอธิบายภาษาคนมีความหมายใกล้เคียงกับคำอธิบายเทคนิคในฐานความรู้มากกว่า
  \item ค่า $k$ มีผลโดยตรง หาก $k$ น้อยเกินไปเฉลยอาจหลุดจาก Top-$k$ (การทดลอง 3.2)
\end{enumerate}

\subsection{ข้อจำกัดของการทดลอง}
การทดลองนี้ใช้เพียง 1 เทคนิคและ 1 กรณี และคำสั่ง \texttt{net users} เป็นคำสั่งที่มีลักษณะชัดเจน
จึงยังไม่สามารถสรุปผลโดยทั่วไปได้ จำเป็นต้องทดสอบกับชุดข้อมูลที่ครอบคลุมหลายเทคนิคเพิ่มเติม
และกำหนดเกณฑ์การประเมินไว้ล่วงหน้าก่อนดูผล

\section{แผนการทดลองหลัก}
\label{sec:main-exp}

การทดลองหลักใช้ Development Set และ Test Set ตามหัวข้อ~\ref{sec:dataset}
โดยปรับพารามิเตอร์และ Prompt บน Development Set แล้วล็อกค่า (Freeze)
ก่อนรันบน Test Set เพียงครั้งเดียว

\subsection{Baseline และหลักการเปลี่ยนทีละปัจจัย}
Baseline ของการทดลองคือ Event Selection แบบ Top-$k$ ที่งบประมาณ $B^{*}$
นำเสนอในรูปแบบ JSON และส่งให้ LLM โดยตรงโดยไม่ใช้ RAG (Direct-Log)
โดย $B^{*}$ คือค่า $B \in \{500, 1{,}000, 2{,}000, 4{,}000\}$ Tokens ที่ได้ความถูกต้องสูงสุดบน Development Set
ค่า $B$ สูงสุดกำหนดที่ 4,000 Tokens เพราะ Context Window ขนาด 16,384 Tokens
ต้องเผื่อพื้นที่ให้ Prompt ผลลัพธ์จาก RAG ($k=10$ ประมาณ $10 \times 280$ Tokens) และคำตอบของโมเดล
และการเพิ่มค่าทีละสองเท่าช่วยให้เห็นจุดที่ความถูกต้องหยุดเพิ่มขึ้น

การทดลองใช้หลักเปลี่ยนทีละปัจจัยจาก Baseline (One-factor-at-a-time)
เนื่องจากการทดลองทุกการผสม ($5 \times 5 \times 5 = 125$ รูปแบบ $\times$ 40 กรณี = 5,000 ครั้ง)
ใช้เวลามากเกินกว่าที่โครงงานทำได้ ข้อจำกัดคือไม่เห็นปฏิสัมพันธ์ระหว่างปัจจัย

\subsection{R0: Raw Logs}
เพื่อตอบวัตถุประสงค์ว่า Investigation Context ดีกว่าการใช้ Raw Logs หรือไม่
R0 ส่ง Alert พร้อม Events ดิบทั้งหมดในช่วงเวลา $\pm T$ รอบ Alert ในรูปแบบต้นฉบับ
โดยไม่คัดเลือก ไม่ Normalize และไม่ตัด Field ออก
หากข้อมูลยาวเกิน Context Window จะลดค่า $T$ จนพอดีและบันทึกค่า $T$ ที่ใช้จริง

\subsection{E1: Event Selection}
เปรียบเทียบวิธี Event Selection ทั้ง 5 วิธี โดยใช้รูปแบบ JSON และไม่ใช้ RAG (ตารางที่~\ref{tab:exp-e1})
วิธีอื่นนอกจาก Top-$k$ ไม่จำกัดด้วยงบประมาณ $B$ แต่บันทึกจำนวน Token ทุกครั้ง

\begin{table}[ht]
  \centering
  \caption{รูปแบบการทดลอง E1: Event Selection}
  \label{tab:exp-e1}
  \small
  \begin{tabular}{|l|l|L{0.5\textwidth}|}
    \hline
    \textbf{รหัส} & \textbf{วิธี} & \textbf{สิ่งที่ส่งให้ LLM} \\ \hline
    E1-1 & Alert Only & เฉพาะ Field ของ Alert \\ \hline
    E1-2 & Single Event & Event ทั้งหมดที่ Alert ชี้ (1 Event) \\ \hline
    E1-3 & Time Window & Events ที่ Normalize แล้วทั้งหมดในช่วง $\pm T$ \\ \hline
    E1-4 & Multiple Related & Events ในสายความสัมพันธ์ ProcessGuid (บรรพบุรุษและลูกหลาน) ไม่ให้คะแนนและไม่จำกัดงบประมาณ \\ \hline
    E1-5 ถึง E1-8 & Top-$k$ & Graph $\rightarrow$ Traverse $\rightarrow$ Score $\rightarrow$ Budget ที่ $B$ = 500, 1,000, 2,000 และ 4,000 \\ \hline
  \end{tabular}
\end{table}

\subsection{E2: Representation}
ใช้ Top-$k$ ที่ $B^{*}$ และไม่ใช้ RAG แล้วเปลี่ยนเฉพาะรูปแบบการนำเสนอข้อมูล (ตารางที่~\ref{tab:exp-e2})
การคัดเลือก Events ใช้งบประมาณที่นับเป็น Token ของ JSON แล้วแปลงชุด Events เดียวกันเป็นรูปแบบอื่น
ทุกรูปแบบจึงมีข้อมูลชุดเดียวกัน ข้อจำกัดคือรูปแบบ Timeline ตัด Field บางส่วนออก

\begin{table}[ht]
  \centering
  \caption{รูปแบบการทดลอง E2: Representation}
  \label{tab:exp-e2}
  \small
  \begin{tabular}{|l|l|c|}
    \hline
    \textbf{รหัส} & \textbf{รูปแบบ} & \textbf{Tokens ต่อ Event (ตัวอย่าง)} \\ \hline
    Baseline & JSON & $\approx$ 138 \\ \hline
    E2-1 & XML & $\approx$ 195 \\ \hline
    E2-2 & TOON & $\approx$ 119 \\ \hline
    E2-3 & Key=Value & $\approx$ 120 \\ \hline
    E2-4 & Timeline & $\approx$ 35 \\ \hline
  \end{tabular}
\end{table}

\subsection{E3: RAG}
ใช้ Top-$k$ ที่ $B^{*}$ และรูปแบบ JSON แล้วเปลี่ยนวิธีการใช้ RAG (ตารางที่~\ref{tab:exp-e3})
ซึ่งขยายจากการทดลองนำร่องในหัวข้อก่อนหน้าให้ครอบคลุมทุกกรณีใน Test Set
ย่อหน้าสรุปพฤติกรรม (Summary) สร้างครั้งเดียวต่อกรณีและใช้ร่วมกันใน E3-1 และ E3-4
โดยนับจำนวน Token และเวลาของการสร้างย่อหน้าสรุปรวมไว้ด้วย

\begin{table}[ht]
  \centering
  \caption{รูปแบบการทดลอง E3: RAG}
  \label{tab:exp-e3}
  \small
  \begin{tabular}{|l|l|L{0.22\textwidth}|L{0.28\textwidth}|}
    \hline
    \textbf{รหัส} & \textbf{ชื่อ} & \textbf{Query ของ RAG} & \textbf{สิ่งที่ LLM เห็น} \\ \hline
    Baseline & Direct-Log & -- & Logs \\ \hline
    E3-1 & Direct-Summary & -- & ย่อหน้าสรุป 80--150 คำ (ไม่มีชื่อหรือรหัสเทคนิค) \\ \hline
    E3-2 & RAG-Log ($k=10$) & Logs JSON ทั้งหมด & Logs + 10 เทคนิคที่ค้นคืนได้ \\ \hline
    E3-3 & RAG-Log ($k=5$) & Logs JSON ทั้งหมด & Logs + 5 เทคนิคที่ค้นคืนได้ \\ \hline
    E3-4 & RAG-Summary ($k=10$) & ย่อหน้าสรุป & Logs + 10 เทคนิคที่ค้นคืนได้ \\ \hline
  \end{tabular}
\end{table}

\noindent
ค่า $k$ ของ RAG คือจำนวนเทคนิคที่ค้นคืนมา ซึ่งเป็นคนละความหมายกับ Top-$k$ ของ Event Selection

\subsection{จำนวนการทดลองและค่าที่ควบคุม}
การทดลองบน Test Set มีทั้งหมด 17 รูปแบบ (ตารางที่~\ref{tab:exp-runs})
ค่าที่คงที่ตลอดการทดลอง ได้แก่ LLM Qwen3.5 9B ผ่าน Ollama, Temperature = 0, ปิดโหมด Thinking,
Context Window 16,384 Tokens, Prompt ชุดเดียวกัน และฐานความรู้ ATT\&CK v19.2
ส่วนค่าที่ปรับบน Development Set แล้วล็อกไว้ ได้แก่ $\gamma$, $\tau$, $w[\mathrm{EventID}]$,
จำนวน Hop สูงสุด, $B^{*}$, ค่า $\pm T$ ของ Time Window และ R0 และ Prompt

\begin{table}[ht]
  \centering
  \caption{จำนวนรูปแบบการทดลองและจำนวนครั้งที่รันบน Test Set}
  \label{tab:exp-runs}
  \small
  \begin{tabular}{|l|c|c|}
    \hline
    \textbf{ชุดการทดลอง} & \textbf{จำนวนรูปแบบ} & \textbf{จำนวนครั้ง (40 กรณี)} \\ \hline
    R0 & 1 & 40 \\ \hline
    E1 (รวม Baseline) & 8 & 320 \\ \hline
    E2 & 4 & 160 \\ \hline
    E3 & 4 & 160 \\ \hline
    \textbf{รวม} & \textbf{17} & \textbf{680 (+ 40 ครั้งสำหรับสร้างย่อหน้าสรุป)} \\ \hline
  \end{tabular}
\end{table}

\subsection{ตัวชี้วัดและการวิเคราะห์ทางสถิติ}
ตัวชี้วัดที่ใช้ตามหัวข้อ~\ref{sec:metrics} จับคู่กับวัตถุประสงค์ของโครงงานดังตารางที่~\ref{tab:exp-metrics}
ความถูกต้องรายงานพร้อมช่วงความเชื่อมั่น 95\% แบบ Bootstrap
และเปรียบเทียบแต่ละรูปแบบกับ Baseline ด้วย McNemar's Test จำนวน 16 คู่ โดยปรับค่า $p$ ด้วยวิธีของ Holm
นอกจากนี้จะรายงานผลแยกตามเทคนิค และแยกกรณีที่ใช้ Anchor สำรองออกจากกรณีที่มี Alert จริง

\begin{table}[ht]
  \centering
  \caption{ตัวชี้วัดที่ใช้ตอบวัตถุประสงค์ของโครงงาน}
  \label{tab:exp-metrics}
  \small
  \begin{tabular}{|L{0.3\textwidth}|L{0.38\textwidth}|l|}
    \hline
    \textbf{วัตถุประสงค์} & \textbf{ตัวชี้วัด} & \textbf{วัดบน} \\ \hline
    ความถูกต้องในการระบุเทคนิค & Exact Match และ Parent Match Accuracy & Test Set \\ \hline
    การลดปริมาณข้อมูล & Input Tokens, Latency และสัดส่วน Events ที่ถูกคัดเลือก & Test Set \\ \hline
    การคงหลักฐานที่จำเป็น & Evidence Recall & Development Set \\ \hline
    ประสิทธิภาพของ RAG (E3) & อันดับของเฉลยและ hit@$k$ & Test Set \\ \hline
    ข้อมูลประกอบปัญหา & จำนวน Alert ต่อไฟล์ และจำนวน TP/FP & ทั้งสองชุด \\ \hline
  \end{tabular}
\end{table}

\section{ลำดับการดำเนินการทดลอง}
\begin{table}[ht]
  \centering
  \caption{ลำดับการดำเนินการทดลอง}
  \label{tab:exp-phases}
  \small
  \begin{tabular}{|c|L{0.55\textwidth}|l|}
    \hline
    \textbf{ระยะ} & \textbf{งาน} & \textbf{ข้อมูล} \\ \hline
    0 & การทดลองนำร่อง (เกือบเสร็จสิ้น) & T1087.001 \\ \hline
    1 & Freeze, Normalize, Screening, สร้าง Alert และกรณีทดสอบ & Dev + Test \\ \hline
    2 & ระบุ Event สำคัญอัตโนมัติจาก Process Tree ของ Atomic Red Team แล้วสุ่มตรวจ & Dev \\ \hline
    3 & ปรับพารามิเตอร์แล้วล็อกค่า & Dev \\ \hline
    4 & รันทั้ง 17 รูปแบบเพียงครั้งเดียว & Test \\ \hline
  \end{tabular}
\end{table}

\section{ประเด็นที่ยังเปิดและข้อจำกัด}
\begin{itemize}
  \item \textbf{การเลือก Alert เมื่อมีหลายรายการ} -- รอบนี้ใช้ 1 Alert ต่อไฟล์ การเลือกหรือจัดกลุ่ม Alert จะศึกษาต่อ
  \item \textbf{ป้ายกำกับที่น่าสงสัยใน OTRF} -- ต้องทบทวนก่อนล็อกค่าการทดลอง
  \item \textbf{ตัวแปลง OTRF JSON เป็น Sysmon XML} -- ยังต้องพัฒนา และเป็นความเสี่ยงหลักของการสร้าง Alert ฝั่ง Test Set
  \item \textbf{จำนวนข้อมูลเป็นค่าสูงสุด} -- ไฟล์ของ Splunk ยังต้องผ่านการคัดกรองบนเครื่องจริง
        และชุดของ OTRF ยังต้องตรวจเทียบกับ Adversary View จำนวนกรณีจริงจึงอาจลดลง
  \item \textbf{นอกขอบเขต (งานในอนาคต)} -- การจำแนก True Positive/False Positive,
        การประเมินบน Linux Auditd, ชุดข้อมูลแบบ Compound หรือ APT
        และชุดข้อมูลของ Splunk ที่ไม่ได้สร้างด้วย Atomic Red Team
\end{itemize}
